\documentclass[10pt]{article}
\usepackage[utf8]{inputenc}
\usepackage{amsmath, amssymb}
\usepackage{graphicx}
\usepackage{wrapfig}
\usepackage{multirow}
\usepackage[table]{xcolor} % Added [table] option to fix the \rowcolor error
\usepackage{listings}
\usepackage{hyperref}
\usepackage{authblk}

\hypersetup{
    colorlinks=true,
    linkcolor=blue,
    filecolor=magenta,      
    urlcolor=cyan,
}

\title{\textbf{Laplace Analysis and Applications}}
\author[1]{Pierre-Simon Laplace}
\author[2]{Oliver Heaviside}
\author[3]{Gustav Robert Kirchhoff}
\affil[1]{Académie des Sciences, Paris, France}
\affil[2]{Royal Society, London, United Kingdom}
\affil[3]{University of Heidelberg, Heidelberg, Germany}
\date{}

\begin{document}

\maketitle

\section{From Functions to Transforms}
\subsection{The Laplace Transform}
Let $f(t)$ be a function defined for $t\ge0$. Its Laplace transform is defined by
\begin{equation}
\mathcal{L}\{f(t)\}=F(s)=\int_{0}^{\infty}e^{-st}f(t)dt, \quad s>0
\end{equation}
The transform changes a function of time into a function of the new variable $s$.\footnote{The variable $s$ is commonly interpreted as a complex-frequency variable.} Symbolically,
\begin{equation}
f(t)\rightarrow F(s)
\end{equation}
The inverse Laplace transform reverses this process:
\begin{equation}
f(t)=\mathcal{L}^{-1}\{F(s)\}.
\end{equation}
The main purpose of the method is to replace differentiation and integration in the time domain by simpler algebraic operations in the transform domain \cite{kreyszig}. A concise overview is available from the \href{https://mathworld.wolfram.com/LaplaceTransform.html}{Laplace Transform} entry on MathWorld.

\subsection{A Simple Example}
For the constant function $f(t)=1$,
\begin{align}
\mathcal{L}\{1\} &= \int_{0}^{\infty}e^{-st}dt \nonumber \\
&= \left[-\frac{1}{s}e^{-st}\right]_{0}^{\infty} \nonumber \\
&= \frac{1}{s}
\end{align}
Similarly,
\begin{equation}
\mathcal{L}\{t\}=\frac{1}{s^{2}}, \quad \mathcal{L}\{t^{2}\}=\frac{2}{s^{3}}
\end{equation}

\section{Common Transform Equations}
Table \ref{tab:transforms} lists several frequently used Laplace-transform equations.

\begin{table}[h!]
\centering
\caption{Selected Laplace-transform equations.}
\label{tab:transforms}
\begin{tabular}{|l|c|c|l|}
\hline
\rowcolor{gray!20}
\textbf{Family} & \textbf{Time function} & \textbf{Transform Equation} & \textbf{Condition} \\ \hline
Laplace transform & $1$ & $\frac{1}{s}$ & $s>0$ \\ \hline
\multirow{2}{*}{Algebraic} & $t$ & $\frac{1}{s^{2}}$ & \multirow{2}{*}{$s>0, n=0,1,2,...$} \\ 
 & $t^{n}$ & $\frac{n!}{s^{n+1}}$ & \\ \hline
\multirow{2}{*}{Exponential} & $e^{at}$ & $\frac{1}{s-a}$ & $s>a$ \\ 
 & $e^{-at}$ & $\frac{1}{s+a}$ & $s>-a$ \\ \hline
\multirow{2}{*}{Oscillatory} & $\sin(\omega t)$ & $\frac{\omega}{s^{2}+\omega^{2}}$ & $\omega>0$ \\ 
 & $\cos(\omega t)$ & $\frac{s}{s^{2}+\omega^{2}}$ & $\omega>0$ \\ \hline
 & Unit-step function $u(t-a)$ & $\frac{e^{-as}}{s}$ & $a>0$ \\ \hline
\end{tabular}
\end{table}

The formulas in Table \ref{tab:transforms} allow many functions to be transformed without repeatedly evaluating improper integrals.

\section{Application to an RC Circuit}
\begin{wrapfigure}{r}{0.4\textwidth}
  \centering
  \includegraphics[width=\linewidth]{rc_circuit.png}
  \caption{A simple series RC circuit.}
  \label{fig:circuit}
\end{wrapfigure}

\textbf{Circuit Model} \\
Figure \ref{fig:circuit} shows a simple resistor--capacitor circuit connected to an input voltage source. The resistor limits the current, while the capacitor stores electrical energy.

For a resistor $R$ and capacitor $C$, Kirchhoff's voltage law gives
\begin{equation}
RC\frac{dv_{C}(t)}{dt}+v_{C}(t)=v_{in}(t),
\end{equation}
where $v_{C}(t)$ is the voltage across the capacitor. Assuming $v_{C}(0)=0$, taking the Laplace transform of Equation (5) gives
\begin{align}
RC[sV_{C}(s)-v_{C}(0)]+V_{C}(s) &= V_{in}(s), \nonumber \\
\Rightarrow RCsV_{C}(s)+V_{C}(s) &= V_{in}(s) \nonumber \\
\Rightarrow V_{C}(s)(RCs+1) &= V_{in}(s).
\end{align}
The transfer function is therefore
\begin{equation}
H(s)=\frac{V_{C}(s)}{V_{in}(s)}=\frac{1}{RCs+1}.
\end{equation}
For a constant input $V_{0}$,
\begin{equation}
V_{in}(s)=\frac{V_{0}}{s}.
\end{equation}
Thus,
\begin{equation}
V_{C}(s)=\frac{V_{0}}{s(RCs+1)}.
\end{equation}
Taking the inverse transform gives
\begin{equation}
v_{C}(t)=V_{0}(1-e^{-t/(RC)}).
\end{equation}
The quantity $\tau=RC$ is called the time constant of the circuit. A larger value of $\tau$ causes the capacitor to charge more slowly \cite{ogata}.

\section{How Laplace Analysis Is Used}
The Laplace transform is useful in many areas:
\begin{itemize}
    \item \textbf{Differential equations}
    \begin{itemize}
        \item initial conditions are included automatically;
        \item derivatives become algebraic expressions.
    \end{itemize}
    \item \textbf{Electrical circuits}
    \begin{itemize}
        \item voltages and currents are studied in the s-domain;
        \item components lead to transfer functions.
    \end{itemize}
    \item \textbf{Control and signal systems}
    \begin{itemize}
        \item stability can be studied through poles;
        \item common test inputs include:
        \begin{itemize}
            \item[$\triangleright$] the unit-step input;
            \item[$\triangleright$] the impulse input;
            \item[$\triangleright$] sinusoidal inputs.
        \end{itemize}
    \end{itemize}
\end{itemize}
In the time domain, a system may be described by a differential equation, whereas in the Laplace domain, the same system is often described by a rational algebraic expression.

\section{A Short Python Implementation}
Listing 1 evaluates the capacitor voltage at a finite set of time instants using an external Python file.

\lstset{language=Python, numbers=left, frame=single, basicstyle=\ttfamily\small, keywordstyle=\color{blue}}
\begin{lstlisting}[caption={Computing samples of the RC step response.}, label={lst:python}]
from math import exp
def rc_response(resistance, capacitance, input_voltage, times):
    """Return capacitor-voltage samples for an RC step input."""
    if resistance <= 0 or capacitance <= 0:
        raise ValueError("R and C must be positive")
    
    tau = resistance * capacitance
    values = []
    
    for time in times:
        voltage = input_voltage * (1 - exp(-time / tau))
        values.append(voltage)
        
    return values

times = [0.0, 0.5, 1.0, 2.0]
print(rc_response(2.0, 0.5, 5.0, times))
\end{lstlisting}

Equation (9) and Listing 1 describe the same computation.

\begin{thebibliography}{9}
\bibitem{kreyszig}
E. Kreyszig, \textit{Advanced Engineering Mathematics}, 10th ed. John Wiley \& Sons, 2011.
\bibitem{ogata}
K. Ogata, \textit{Modern Control Engineering}, 5th ed. Prentice Hall, 2010.
\end{thebibliography}

\end{document}